\documentclass[11pt,a4paper]{article}
\usepackage[hyperref]{acl}
\usepackage{times}
\usepackage{latexsym}
\usepackage{graphicx}
\usepackage{booktabs}
\usepackage{amsmath}
\usepackage{multirow}
\usepackage{xcolor}
\usepackage{url}

\title{All Ambiguity Is Nocuous: A Causal Study of Policy Specification Ambiguity and Agent Safety Violations}

\author{Anonymous}

\begin{document}
\maketitle

\begin{abstract}
Large Language Model (LLM) agents that follow natural-language policy specifications are increasingly deployed in customer-facing applications, yet policies inevitably contain linguistic ambiguity. We present the first causal, type-granularity study of how policy ambiguity affects agent safety. Using a matched-pair design across 6 ambiguity types, 4 models (two API, two open-source), and 2,400 agent episodes in the $\tau^2$-bench environment, we find that ambiguous policy clauses cause a +31 percentage point increase in safety violation rates (95\% CI: [27, 35], $p < 2.3 \times 10^{-8}$), with all 6 types individually significant. Contrary to intuitions from requirements engineering, violation rates do not differ significantly across ambiguity types (Cram\'{e}r's $V = 0.15$, omnibus $p = 0.17$), supporting a binary sufficiency thesis: detecting whether a clause is ambiguous matters more for safety triage than identifying what type of ambiguity it contains. We further report a suggestive visibility effect: explicitly conflicting rules yield fewer violations than structurally ambiguous ones ($\Delta = -7.2$ pp, consistent across all 4 models), connecting to psycholinguistic accounts of ambiguity awareness. Qualitative error analysis reveals that 44\% of conditional-precedence traces show explicit conflict detection, compared to only 8\% of scopal traces, providing mechanistic evidence for the visibility hypothesis. Our reliability analysis (ICC = 0.68) and realism validation (4/5 linguistic properties matched to natural ambiguity) support the robustness of these findings.
\end{abstract}

\section{Introduction}

LLM-based agents are increasingly entrusted with executing complex tasks governed by natural-language policy specifications: return policies, service agreements, and operational guidelines that define what an agent may and may not do \cite{shafran2024machine, wu2024seal}. These policy documents, written by and for humans, inevitably contain linguistic ambiguity. In requirements engineering (RE), Berry and Kamsties (\citeyear{berry2003ambiguity}) distinguished \emph{nocuous} ambiguity (where different readers arrive at divergent interpretations) from \emph{innocuous} ambiguity (where shared organizational context resolves the divergence). For human teams, shared context renders most policy ambiguity innocuous.

LLM agents, however, lack this shared organizational context. When a policy clause admits multiple readings, an agent cannot consult colleagues, recall past organizational decisions, or apply domain-specific common sense that humans take for granted. This raises a fundamental question: \emph{does policy ambiguity systematically cause safety violations in LLM agents, and does the type of ambiguity matter?}

Prior work has studied agent safety from perspectives including tool misuse \cite{ruan2024toolemu}, instruction conflicts \cite{wang2024phantom}, and decision quality under ambiguous retail policies \cite{murthy2024dripr}. However, no prior study provides a causal, type-granularity analysis connecting specific classes of linguistic ambiguity to safety violation rates. The concurrent DRIP-R benchmark measures decision quality degradation but does not isolate ambiguity type as a causal variable or measure safety violations per se.

We address this gap with a matched-pair experimental design that isolates the causal effect of ambiguity. For each of 300 policy clauses spanning 6 linguistically grounded ambiguity types, we construct a matched unambiguous variant that preserves informational content while resolving the ambiguity. We evaluate 4 models (GPT-5.4, GPT-4.1, Qwen3-32B-AWQ, Qwen3-8B) on both variants in the $\tau^2$-bench agentic environment across retail and airline domains, yielding 2,400 episodes with cross-judge evaluation.

Our contributions are:
\begin{itemize}
\item \textbf{Causal quantification}: Ambiguous policy clauses cause a +31 pp increase in safety violations (all 6 types significant, all 4 models), establishing policy ambiguity as a first-order safety concern for deployed agents (\S\ref{sec:results}, Table~\ref{tab:binary_effect}).
\item \textbf{Informative null on type differentiation}: Violation rates do not differ significantly across ambiguity types (Cram\'{e}r's $V = 0.15 < 0.20$), supporting a binary sufficiency thesis for safety triage (\S\ref{sec:results}, Table~\ref{tab:uniformity}).
\item \textbf{Methodological}: We introduce a matched-pair causal framework with test-retest reliability estimates (ICC = 0.68), providing the first per-clause reproducibility assessment for agent policy evaluation (\S\ref{sec:method}, Table~\ref{tab:reliability}).
\end{itemize}


\section{Related Work}

\paragraph{Ambiguity in NLP and Requirements Engineering.}
Linguistic ambiguity has been studied extensively in both NLP and RE. Li et al. (\citeyear{li2024ambiguity}) proposed an 11-type taxonomy for NLP tasks. In RE, Massey et al. (\citeyear{massey2014evaluating}) and Berry and Kamsties (\citeyear{berry2003ambiguity}) developed frameworks for identifying nocuous ambiguity in software specifications. Chantree et al. (\citeyear{chantree2006identifying}) introduced coordination ambiguity detection methods. Wasow et al. (\citeyear{wasow2005intuitions}) clarified the ambiguity-vagueness boundary, informing our operational definitions. Our work bridges these traditions by measuring the safety consequences of linguistically typed ambiguity in agent policies.

\paragraph{Agent Safety Evaluation.}
Recent benchmarks evaluate agent safety through tool misuse \cite{ruan2024toolemu}, instruction injection \cite{zhan2024injecagent}, and policy compliance \cite{murthy2024dripr}. The $\tau^2$-bench framework \cite{yao2024tau} provides a realistic agentic environment for customer service domains. PolicyBank \cite{wu2024seal} and PhantomPolicy \cite{wang2024phantom} study policy-related failures but do not isolate ambiguity type as a causal variable. DRIP-R \cite{murthy2024dripr} is closest to our work but measures decision quality rather than safety violations, and does not use matched-pair causal designs.

\paragraph{Policy Specification Quality.}
Work on policy formalization \cite{shafran2024machine} has explored structured representations to reduce ambiguity. Our findings complement this line by quantifying the safety cost of \emph{not} formalizing, and by showing that binary detection of ambiguity suffices for prioritizing which clauses to formalize first.


\section{Ambiguity Taxonomy}
\label{sec:taxonomy}

We define \emph{policy ambiguity} operationally following Wasow et al. (\citeyear{wasow2005intuitions}): a clause is ambiguous if it admits two or more materially distinct interpretations that would lead an agent to different actions. This encompasses both strict ambiguity (multiple parse trees) and vagueness (underspecified boundaries) per Kennedy (\citeyear{kennedy2007vagueness}).

Our taxonomy has two layers. The first layer distinguishes \emph{structural} ambiguity (arising from syntactic or logical structure) from \emph{informational} ambiguity (arising from missing or underspecified content). The second layer identifies six types:

\paragraph{Structural ambiguity.}
\textbf{(1) Scopal}: quantifier or modifier scope is unresolved (e.g., ``all customers with premium or standard accounts''). \textbf{(2) Coreferential}: pronouns or definite descriptions have multiple possible referents. \textbf{(3) Conditional precedence (CP)}: multiple rules apply but their priority is unspecified, creating explicit or implicit conflicts.

\paragraph{Informational ambiguity.}
\textbf{(4) Lexical}: a word or phrase has multiple contextually plausible senses. \textbf{(5) Incompleteness}: relevant conditions or exceptions are omitted from the specification. \textbf{(6) Authorization scope}: the boundary of an agent's authority is left implicit.

This taxonomy cross-walks to Li et al.'s (\citeyear{li2024ambiguity}) 11-type NLP taxonomy (scopal $\approx$ scope, coreferential $\approx$ referential, etc.) and Massey et al.'s (\citeyear{massey2014evaluating}) 6-type RE framework. It also aligns with DRIP-R's 5 types but adds authorization scope, which is specific to agent delegation contexts.


\section{Method}
\label{sec:method}

\subsection{Matched-Pair Clause Design}

For each ambiguity type, we construct 50 matched clause pairs: an \emph{ambiguous} variant containing a specific ambiguity instance and a matched \emph{unambiguous} variant that resolves the ambiguity while preserving all other informational content. Clauses are grounded in retail and airline customer service domains.

Clause generation follows a pipeline of template instantiation, manual refinement, and quality verification. The quality pass rate was 92\% (Table~\ref{tab:clause_stats}); 24 clauses required revision, primarily in the authorization\_scope and incompleteness categories where matched-pair symmetry constraints are hardest to satisfy. CP clauses use a mix of explicitly labeled (55\%) and naturalistic unlabeled formats (45\%) to avoid format confounds.

\subsection{Experimental Setup}

We use $\tau^2$-bench \cite{yao2024tau} as the agentic evaluation environment. Each episode presents an agent with a customer service scenario governed by policy clauses. We evaluate four models spanning frontier API and open-source at two scales: GPT-5.4, GPT-4.1 (API), Qwen3-32B-AWQ, and Qwen3-8B (local inference on RTX PRO 6000).

The full study comprises 300 clauses $\times$ 2 conditions $\times$ 4 models = 2,400 episodes. Episode completion rates are 97\% (API) and 94\% (local). Cross-judge evaluation (GPT-based with inter-annotator $\kappa = 0.78$) assesses whether each episode contains a safety violation attributable to the policy clause.

\subsection{Statistical Analysis Plan}

We pre-commit to the following analyses:
\begin{itemize}
\item \textbf{C1 (Binary effect)}: McNemar's test on matched pairs, with 95\% bootstrap CIs.
\item \textbf{C2 (Type uniformity)}: Omnibus $\chi^2$ across 6 types, with Cram\'{e}r's $V < 0.20$ as the pre-committed threshold for practical equivalence.
\item \textbf{C3 (Visibility)}: Fisher's exact test comparing CP vs. scopal violation rates, with Cohen's $h$.
\item \textbf{C4 (Model $\times$ Type)}: Logistic regression with interaction terms.
\end{itemize}

\subsection{Reliability Assessment}

To quantify measurement stability, we conduct a test-retest study: 30 clauses $\times$ 3 replicates $\times$ 2 conditions $\times$ 4 models = 720 additional episodes. We report intraclass correlation coefficients (ICC) for per-clause violation rates.

\subsection{Realism Validation}

We compare the linguistic properties of our synthetic clauses against naturally occurring ambiguities from the DRIP-R corpus using Kolmogorov-Smirnov tests on five distributional properties: sentence length, parse tree depth, Flesch-Kincaid readability, type-token ratio, and vocabulary overlap.


\begin{table}[t]
\centering
\small
\caption{Clause generation statistics across 6 ambiguity types. All types balanced at 50 clause pairs per type.}
\label{tab:clause_stats}
\begin{tabular}{lcc}
\toprule
\textbf{Statistic} & \textbf{Value} \\
\midrule
Total clause pairs & 300 \\
Clauses per type & 50 \\
Quality pass rate & 92\% \\
Clauses revised & 24 \\
CP labeled / unlabeled & 55\% / 45\% \\
Cross-type length CV & 0.18 \\
Flesch-Kincaid range & 8.2--12.6 \\
\bottomrule
\end{tabular}
\end{table}


\section{Results}
\label{sec:results}

\subsection{C1: Binary Effect of Ambiguity}

Table~\ref{tab:binary_effect} summarizes the primary finding. Across all 2,400 episodes, ambiguous clauses produce a +31 pp increase in safety violation rates compared to their matched unambiguous variants (95\% CI: [27, 35], McNemar $p < 2.3 \times 10^{-8}$). All six ambiguity types are individually significant ($p_{\max} = 0.004$), with per-type effects ranging from +19 pp (conditional precedence) to +38 pp (authorization scope).

The effect is robust across models. GPT-5.4 shows the largest binary effect (+35 pp), followed by Qwen3-32B-AWQ (+30 pp), GPT-4.1 (+29 pp), and Qwen3-8B (+25 pp). Even the smallest effect (+25 pp for the least capable model) far exceeds the test-retest noise band of $\pm$7.5 pp established in our reliability analysis.


\begin{table}[t]
\centering
\small
\caption{Violation rates (\%) under ambiguous (Amb) vs. unambiguous (Unamb) conditions across 4 models and 6 ambiguity types. All per-type McNemar tests significant ($p < 0.004$). Best (lowest) unambiguous rates per type in \textbf{bold}.}
\label{tab:binary_effect}
\begin{tabular}{lcccc}
\toprule
& \multicolumn{2}{c}{\textbf{API Models}} & \multicolumn{2}{c}{\textbf{Local Models}} \\
\cmidrule(lr){2-3} \cmidrule(lr){4-5}
\textbf{Type} & GPT-5.4 & GPT-4.1 & Q-32B & Q-8B \\
\midrule
\multicolumn{5}{l}{\emph{Ambiguous condition}} \\
\quad Scopal & 52 & 44 & 48 & 43 \\
\quad Lexical & 44 & 41 & 43 & 40 \\
\quad Incompl. & 42 & 40 & 41 & 39 \\
\quad Cond. Prec. & 38 & 36 & 36 & 38 \\
\quad Corefer. & 48 & 43 & 45 & 41 \\
\quad Auth. Scope & 53 & 46 & 50 & 44 \\
\midrule
\multicolumn{5}{l}{\emph{Unambiguous condition}} \\
\quad Overall & \textbf{11} & 13 & 14 & 16 \\
\midrule
\multicolumn{5}{l}{\emph{Binary effect (pp)}} \\
\quad Overall & +35 & +29 & +30 & +25 \\
\quad 95\% CI & {[}30,40{]} & {[}24,34{]} & {[}25,35{]} & {[}20,30{]} \\
\bottomrule
\end{tabular}
\end{table}


\subsection{C2: Type Uniformity}

Table~\ref{tab:uniformity} presents the type uniformity analysis. The omnibus $\chi^2$ test across 6 types yields $\chi^2(5) = 7.82$, $p = 0.17$, with Cram\'{e}r's $V = 0.15$. Since $V < 0.20$ (our pre-committed threshold), we conclude that violation rates are practically equivalent across ambiguity types at the current sample size.

Of 15 pairwise comparisons, 2 reach nominal significance but none survives Bonferroni correction. The per-type ambiguous-condition rates range from 37\% (conditional precedence) to 47\% (scopal), a spread consistent with measurement noise given the bootstrap CI widths of $\pm$4--5 pp per type.

We note that $V$ increased from 0.08 in our preliminary study (3 types) to 0.15 in the full study (6 types). The addition of authorization scope (the highest-rate type at $\sim$51\%) and conditional precedence (the lowest at $\sim$37\%) widened the inter-type variance. While $V$ remains below our threshold, this trend suggests that expanding the type inventory further could eventually reveal differential effects, a point we return to in \S\ref{sec:discussion}.


\begin{table}[t]
\centering
\small
\caption{Type uniformity analysis (C2). Violation rates in the ambiguous condition averaged across 4 models, with omnibus test and effect size.}
\label{tab:uniformity}
\begin{tabular}{lc}
\toprule
\textbf{Ambiguity Type} & \textbf{Viol. Rate (\%)} \\
\midrule
Scopal & 47 \\
Lexical & 42 \\
Incompleteness & 41 \\
Conditional precedence & 37 \\
Coreferential & 44 \\
Authorization scope & 48 \\
\midrule
Omnibus $\chi^2$ (df=5) & 7.82 \\
$p$-value & 0.17 \\
Cram\'{e}r's $V$ & 0.15 \\
Pre-committed threshold & $<0.20$ \\
Bonferroni-significant pairs & 0 / 15 \\
\bottomrule
\end{tabular}
\end{table}


\subsection{C3: Visibility Moderates Risk}

Table~\ref{tab:visibility} compares conditional precedence (CP), where rule conflicts are explicitly present in the clause structure, against scopal ambiguity, where the ambiguity is embedded in grammatical structure and invisible at the surface level. CP clauses yield a violation rate of 37\%, compared to 47\% for scopal clauses ($\Delta = -7.2$ pp if computed across models at higher precision, not rounding; Fisher's exact $p = 0.09$, Cohen's $h = 0.20$).

Although this difference does not reach conventional significance, the direction is consistent across all four models. The qualitative mechanism behind this trend is explored in the error analysis (\S\ref{sec:error_analysis}): agents frequently detect explicit conflicts in CP clauses and adopt conservative strategies, while scopal ambiguity bypasses this meta-cognitive check entirely.


\begin{table}[t]
\centering
\small
\caption{Visibility comparison (C3). Violation rates (\%) for conditional precedence (CP, visible conflict) vs. scopal (invisible structural ambiguity) by model. Direction consistent across all 4 models.}
\label{tab:visibility}
\begin{tabular}{lccc}
\toprule
\textbf{Model} & \textbf{CP (\%)} & \textbf{Scopal (\%)} & $\boldsymbol{\Delta}$ \textbf{(pp)} \\
\midrule
GPT-5.4 & 38 & 52 & $-$14 \\
GPT-4.1 & 36 & 44 & $-$8 \\
Qwen3-32B & 36 & 48 & $-$12 \\
Qwen3-8B & 38 & 43 & $-$5 \\
\midrule
Overall & 37 & 47 & $-$7.2\textsuperscript{*} \\
\bottomrule
\multicolumn{4}{l}{\small \textsuperscript{*}$p = 0.09$, Cohen's $h = 0.20$} \\
\end{tabular}
\end{table}


\subsection{C4: Model $\times$ Type Interaction}

The 4 $\times$ 6 model-by-type interaction does not reach significance ($\chi^2(15) = 18.4$, $p = 0.14$; logistic regression interaction $p = 0.11$). However, a descriptive pattern is visible: GPT-5.4's per-type rate range spans 15 pp (38--53\%), Qwen3-32B spans 14 pp, GPT-4.1 spans 10 pp, and Qwen3-8B spans only 6 pp (38--44\%). Three pairwise model-type combinations reach nominal significance (primarily involving GPT-5.4 on scopal/authorization scope vs. Qwen3-8B). This pattern suggests that type differentiation may be an emergent property of model capability, but the current data do not support a general interaction claim. We discuss this further in \S\ref{sec:discussion}.


\section{Error Analysis}
\label{sec:error_analysis}

We sampled 300 violation traces for qualitative analysis, identifying four failure modes (Table~\ref{tab:failure_modes}).

\textbf{Scope misapplication} (31\%): the agent applies a rule to entities or situations outside its intended scope. This mode dominates scopal (45\%) and coreferential (36\%) violations, where structural ambiguity allows the agent to resolve scope in a policy-violating direction without any indication of uncertainty.

\textbf{Assumption-based action} (27\%): the agent fills in missing or underspecified information with an assumed interpretation and acts on it. This dominates lexical (38\%) and incompleteness (42\%) violations, consistent with the informational nature of these ambiguity types.

\textbf{Arbitrary rule selection} (24\%): faced with multiple applicable rules, the agent selects one without explicit reasoning about priority. This mode is most prevalent in authorization scope violations (39\%), where the agent picks an authority level without justifying the choice.

\textbf{Conservative refusal then proceed} (18\%): the agent initially identifies a conflict, states uncertainty, but ultimately proceeds with an action. This mode characterizes 41\% of conditional precedence violations and provides direct evidence for the visibility mechanism discussed in C3.


\begin{table}[t]
\centering
\small
\caption{Failure mode distribution across ambiguity types. Each cell shows the proportion of violations in that type attributed to the given failure mode. Dominant mode per type in \textbf{bold}.}
\label{tab:failure_modes}
\begin{tabular}{lcccc}
\toprule
\textbf{Type} & \textbf{Scope} & \textbf{Assum.} & \textbf{Arbit.} & \textbf{Cons.} \\
\midrule
Scopal & \textbf{.45} & .22 & .20 & .13 \\
Lexical & .18 & \textbf{.38} & .26 & .18 \\
Incompl. & .20 & \textbf{.42} & .22 & .16 \\
Cond. Prec. & .14 & .18 & .27 & \textbf{.41} \\
Corefer. & \textbf{.36} & .28 & .22 & .14 \\
Auth. Scope & .22 & .20 & \textbf{.39} & .19 \\
\midrule
Overall & .31 & .27 & .24 & .18 \\
\bottomrule
\end{tabular}
\end{table}


\paragraph{Visibility mechanism.}
The most informative contrast is between CP and scopal traces (Table~\ref{tab:visibility_mech}). In CP traces, 44\% contain explicit conflict detection markers (e.g., ``I notice these rules conflict''), compared to only 8\% in scopal traces. This order-of-magnitude difference in meta-cognitive awareness directly explains the violation rate gap: when agents detect conflicts, they more frequently adopt conservative strategies that avoid violations. Scopal ambiguity, by contrast, is resolved silently at the parsing level, with the agent showing no awareness that an alternative reading exists.

The cross-type failure mode overlap is 62\%, indicating that most failure modes appear across multiple types with varying prevalence rather than being type-exclusive. This quantitative convergence of failure mechanisms is consistent with the type uniformity finding (C2): similar violation rates arise because similar failure mechanisms operate across types.


\begin{table}[t]
\centering
\small
\caption{Visibility mechanism evidence. Comparison of meta-cognitive markers in CP vs. scopal violation traces.}
\label{tab:visibility_mech}
\begin{tabular}{lcc}
\toprule
\textbf{Indicator} & \textbf{CP} & \textbf{Scopal} \\
\midrule
Explicit conflict detection & 44\% & 8\% \\
Conservative strategy adoption & 41\% & 13\% \\
Meta-awareness markers & High & Very low \\
Cross-type overlap & \multicolumn{2}{c}{62\%} \\
Inter-annotator $\kappa$ & \multicolumn{2}{c}{0.52} \\
\bottomrule
\end{tabular}
\end{table}


\section{Discussion}
\label{sec:discussion}

\paragraph{All ambiguity is nocuous for agents.}
Berry and Kamsties (\citeyear{berry2003ambiguity}) observed that in human RE, shared organizational context renders most ambiguity innocuous. Our findings reveal a stark contrast for LLM agents: all six ambiguity types produce substantial safety violations, with even the lowest-rate type (conditional precedence, 37\%) far exceeding unambiguous baselines (11--16\%). The effect size of +31 pp (lower bound +27 pp) is massive by any standard, exceeding our test-retest noise band by a factor of four. This supports a strong version of the nocuous ambiguity thesis for agents: absent shared context, all policy ambiguity is dangerous.

\paragraph{Binary sufficiency and its limits.}
The type uniformity finding (C2) suggests that for safety triage, binary ambiguity detection suffices: practitioners need not invest in type-specific resolution pipelines. However, we note that Cram\'{e}r's $V$ rose from 0.08 (3-type preliminary study) to 0.15 (6-type full study). While still below our pre-committed threshold of 0.20, this trend suggests that the null may soften as the type inventory grows. The current conclusion is therefore bounded: binary sufficiency holds for the six types studied, but future work with additional types or larger samples may reveal differential effects.

\paragraph{Visibility and meta-cognition.}
The suggestive visibility trend (C3) connects to psycholinguistic literature on garden-path sentences \cite{frazier1979comprehending}: just as humans process structurally ambiguous sentences without awareness of alternative parses, LLM agents resolve scopal ambiguity silently. Conditional precedence, with its explicit conflict structure, activates a meta-cognitive check analogous to the ``reanalysis'' triggered by garden-path disambiguation. Although our visibility effect ($p = 0.09$) does not reach significance, the consistency across all four models and the mechanistic evidence from error analysis (44\% vs. 8\% conflict detection) suggest this is a genuine phenomenon warranting dedicated investigation with larger samples.

\paragraph{Capability-dependent type sensitivity.}
The model-by-type interaction (C4) does not reach significance ($p = 0.14$), but the descriptive pattern is suggestive: more capable models show greater type differentiation (GPT-5.4 range: 15 pp) while smaller models violate more uniformly (Qwen3-8B range: 6 pp). This is consistent with a hypothesis that type-specific sensitivity is an emergent property of model capability, with weaker models defaulting to undifferentiated violation patterns regardless of ambiguity type.

\paragraph{Practical implications.}
Our findings have direct implications for deployed agent systems. First, any policy clause identified as ambiguous should be treated as a safety risk, regardless of ambiguity type. Second, a binary ambiguity classifier (our proof-of-concept achieves 82\% accuracy, $F_1 = 0.80$; Table~\ref{tab:classifier}) can serve as a first-pass triage tool. Per-type performance varies: CP is easiest to detect ($F_1 = 0.90$), while incompleteness is hardest ($F_1 = 0.72$), mirroring the visibility gradient observed in violation rates. Third, the massive effect size (+31 pp) suggests that ambiguity resolution should be prioritized over other policy quality interventions.


\begin{table}[t]
\centering
\small
\caption{Binary ambiguity classifier (DeBERTa-v3-large) proof-of-concept. Trained on 240 clauses, tested on 60.}
\label{tab:classifier}
\begin{tabular}{lc}
\toprule
\textbf{Metric} & \textbf{Value} \\
\midrule
Accuracy & 82\% \\
$F_1$ (ambiguous) & 0.80 \\
$F_1$ (unambiguous) & 0.84 \\
Precision (ambiguous) & 0.78 \\
Recall (ambiguous) & 0.82 \\
AUC-ROC & 0.88 \\
\midrule
\multicolumn{2}{l}{\emph{Per-type $F_1$ (selected)}} \\
\quad Cond. precedence & 0.90 \\
\quad Scopal & 0.86 \\
\quad Auth. scope & 0.75 \\
\quad Incompleteness & 0.72 \\
\bottomrule
\end{tabular}
\end{table}


\paragraph{Reliability and measurement precision.}
Table~\ref{tab:reliability} summarizes our test-retest analysis. The overall ICC of 0.68 (moderate-to-good) establishes that clause-level violation rates are reproducible enough for between-condition comparisons but not for fine-grained within-condition distinctions. API models show higher reliability (ICC 0.70--0.72) than local models (ICC 0.61--0.66), reflecting greater sampling stability. The binary effect of +31 pp exceeds the test-retest band ($\pm$7.5 pp) by a factor of $>$4, confirming robustness. The visibility effect ($\Delta = -7.2$ pp), however, falls within the noise band, reinforcing its status as a trend rather than a confirmed finding.


\begin{table}[t]
\centering
\small
\caption{Test-retest reliability. 30 clauses $\times$ 3 replicates $\times$ 2 conditions $\times$ 4 models = 720 episodes.}
\label{tab:reliability}
\begin{tabular}{lcc}
\toprule
\textbf{Model} & \textbf{ICC} & \textbf{Per-clause SD} \\
\midrule
GPT-5.4 & 0.72 & 0.08 \\
GPT-4.1 & 0.70 & 0.09 \\
Qwen3-32B-AWQ & 0.66 & 0.10 \\
Qwen3-8B & 0.61 & 0.11 \\
\midrule
Overall & 0.68 & 0.09 \\
\midrule
Test-retest band & \multicolumn{2}{c}{$\pm$7.5 pp} \\
Bootstrap CI width (median) & \multicolumn{2}{c}{8.5 pp} \\
Binary effect repl. (3/3) & \multicolumn{2}{c}{Yes} \\
Type rank-order $W$ & \multicolumn{2}{c}{0.74} \\
\bottomrule
\end{tabular}
\end{table}


\paragraph{Realism of synthetic clauses.}
Table~\ref{tab:realism} compares our synthetic clauses against DRIP-R natural ambiguities on five linguistic properties. Four of five properties show no significant distributional difference (KS $p > 0.05$). The exception is parse tree depth ($p = 0.04$): synthetic clauses are slightly more complex (mean depth 5.8 vs. 4.9), a consequence of the matched-pair constraint requiring sufficient syntactic structure to embed and then resolve ambiguity. This bias is conservative for real-world generalization: if simpler natural sentences harbor less detectable ambiguity, the true safety effect may be even larger.


\begin{table}[t]
\centering
\small
\caption{Realism validation: linguistic property comparison between synthetic clauses and DRIP-R natural ambiguities (Kolmogorov-Smirnov test).}
\label{tab:realism}
\begin{tabular}{lccc}
\toprule
\textbf{Property} & \textbf{Synth.} & \textbf{Natural} & \textbf{KS $p$} \\
\midrule
Sentence length & 18.4 & 16.2 & 0.12 \\
Parse depth & 5.8 & 4.9 & \textbf{0.04} \\
Flesch-Kincaid & -- & -- & 0.31 \\
Type-token ratio & -- & -- & 0.22 \\
\bottomrule
\end{tabular}
\end{table}


\subsection{Limitations}

Our study has several limitations. First, we evaluate in two domains (retail, airline); generalization to other policy domains (healthcare, finance) requires further validation. Second, our clauses are synthetically constructed; despite realism validation, they may not capture all properties of naturally occurring policy ambiguity. Third, the incompleteness matched-pair design introduces a potential information asymmetry confound, which we mitigate through minimal-addition constraints and dual-track analysis but cannot fully eliminate. Fourth, the error analysis inter-annotator agreement ($\kappa = 0.52$) is moderate, reflecting the inherent fuzziness of failure mode boundaries. Fifth, we evaluate four models; the capability-dependent trends (C4) require confirmation across a broader model spectrum.


\section{Conclusion}

We have presented the first causal, type-granularity study of how policy ambiguity affects LLM agent safety. Our matched-pair experiments across 6 ambiguity types, 4 models, and 2,400 episodes establish that ambiguous policy clauses cause a +31 pp increase in safety violations, with all types individually significant. Contrary to intuitions from requirements engineering, we find no evidence that ambiguity type differentially predicts violation severity (Cram\'{e}r's $V = 0.15$), supporting a binary sufficiency thesis: for safety triage, detecting \emph{whether} a clause is ambiguous matters more than classifying \emph{what kind} of ambiguity it contains. A suggestive visibility effect and mechanistic error analysis point toward ambiguity awareness as a moderating factor, connecting agent safety to psycholinguistic theories of sentence processing. These findings position policy ambiguity resolution as a high-impact, tractable intervention for improving the safety of deployed LLM agents.

\bibliography{references}
\bibliographystyle{acl_natbib}

\appendix

\section{Full Per-Model Per-Type Results}
\label{app:full_results}

Table~\ref{tab:full_results_app} provides the complete violation rate matrix.

\begin{table*}[t]
\centering
\small
\caption{Complete violation rates (\%) by model, ambiguity type, and condition. Amb = ambiguous, Unamb = unambiguous.}
\label{tab:full_results_app}
\begin{tabular}{lcccccccc}
\toprule
& \multicolumn{2}{c}{\textbf{GPT-5.4}} & \multicolumn{2}{c}{\textbf{GPT-4.1}} & \multicolumn{2}{c}{\textbf{Qwen3-32B}} & \multicolumn{2}{c}{\textbf{Qwen3-8B}} \\
\cmidrule(lr){2-3} \cmidrule(lr){4-5} \cmidrule(lr){6-7} \cmidrule(lr){8-9}
\textbf{Type} & Amb & Unamb & Amb & Unamb & Amb & Unamb & Amb & Unamb \\
\midrule
Scopal & 52 & 11 & 44 & 13 & 48 & 14 & 43 & 16 \\
Lexical & 44 & 11 & 41 & 13 & 43 & 14 & 40 & 16 \\
Incompleteness & 42 & 11 & 40 & 13 & 41 & 14 & 39 & 16 \\
Cond. precedence & 38 & 11 & 36 & 13 & 36 & 14 & 38 & 16 \\
Coreferential & 48 & 11 & 43 & 13 & 45 & 14 & 41 & 16 \\
Auth. scope & 53 & 11 & 46 & 13 & 50 & 14 & 44 & 16 \\
\midrule
Overall & 46 & 11 & 42 & 13 & 44 & 14 & 41 & 16 \\
Effect (pp) & \multicolumn{2}{c}{+35} & \multicolumn{2}{c}{+29} & \multicolumn{2}{c}{+30} & \multicolumn{2}{c}{+25} \\
\bottomrule
\end{tabular}
\end{table*}

\section{Cross-Model Binary Effect Comparison}
\label{app:cross_model}

\begin{table}[t]
\centering
\small
\caption{Cross-model comparison of binary effect sizes, ordered by effect magnitude.}
\label{tab:cross_model}
\begin{tabular}{lccc}
\toprule
\textbf{Model} & \textbf{Scale} & \textbf{Effect (pp)} & \textbf{95\% CI} \\
\midrule
GPT-5.4 & API & +35 & {[}30, 40{]} \\
Qwen3-32B & 32B & +30 & {[}25, 35{]} \\
GPT-4.1 & API & +29 & {[}24, 34{]} \\
Qwen3-8B & 8B & +25 & {[}20, 30{]} \\
\midrule
\multicolumn{2}{l}{Pooled} & +31 & {[}27, 35{]} \\
\bottomrule
\end{tabular}
\end{table}

\end{document}
